\newcommand{\DoNotLoadEpstopdf}{}
\documentclass[a4paper]{scrartcl}
\usepackage{iftex}
\ifPDFTeX
	\usepackage[T1]{fontenc}
	\usepackage{lmodern}
\else
	\usepackage{fontspec}
\fi
\usepackage[ngerman]{babel}
\usepackage{rotating}
\usepackage{geometry}
\usepackage[table]{xcolor}
\usepackage{mathtools}
\usepackage{caption}
\usepackage{amssymb}
\usepackage{pifont}
\usepackage{amsthm}
\usepackage{graphicx}
\usepackage{varwidth}
\usepackage{enumitem}
\usepackage{soul}
\usepackage{array}
\usepackage{zref-lastpage}
\usepackage{wasysym}
\usepackage{wrapfig2}
\usepackage[headsepline=1pt]{scrlayer-scrpage}
\usepackage{xstring}
\usepackage{xspace}
\usepackage{extarrows}
\usepackage[export]{adjustbox}
\usepackage{ifthen}
\usepackage{multicol}
\usepackage{multirow}
\usepackage{framed}
\usepackage{lineno}
\usepackage{siunitx}
\usepackage[thicklines]{cancel}
\usepackage{vectorlogos}
\usepackage{tasks}\DeclareInstance{tasks}{alphabetize}{default}{ label = (\alph*),label-offset=1em}
\settasks{label-width=1.5em}
\usepackage{nicefrac}
\usepackage{changepage}
\usepackage{tabularray}
\usepackage{setspace}
\usepackage{tcolorbox}\tcbuselibrary{skins,raster}
\usepackage{makecell}
\usepackage{interval}\intervalconfig{separator symbol =;}
\usepackage{tikz}
\usetikzlibrary{calc,tikzmark,decorations.text,shapes,patterns.meta,fit,shadows.blur}
\usepackage{tikz-qtree}
\usepackage{tikz-3dplot}
\usepackage{tkz-fct}\tkzfctset{tan style/.style={-}}
\usepackage{tkz-euclide}
\usepackage{colonequals}
\usepackage{fancyqr}
\usepackage{etoolbox}
\usepackage{marginnote}
\usepackage{relsize}
\usepackage{calculator}
\usepackage[hidelinks]{hyperref}

\geometry{left=2cm,right=2cm,top=2cm,bottom=2cm}
\renewcommand{\familydefault}{\sfdefault}
\parindent 0pt
\pagestyle{scrheadings}
\renewcommand*{\headfont}{\normalfont}
\cfoot{}

\theoremstyle{plain}
\newtheorem{Satz}{Satz}
\theoremstyle{definition}
\newtheorem{Aufgabe}{Aufgabe}
\newtheorem{Definition}{Definition}
\let\origgrqq\grqq
\renewcommand{\grqq}{\origgrqq\xspace}
\let\origflqq\flqq
\renewcommand{\flqq}{\origflqq\xspace}
\newcommand{\N}{\mathbb{N}}
\newcommand{\Z}{\mathbb{Z}}
\newcommand{\Q}{\mathbb{Q}}
\newcommand{\R}{\mathbb{R}}
\newcommand{\blau}[1]{\textcolor{blue}{#1}}
\newcommand{\pink}[1]{\textcolor{magenta}{#1}}
\newcommand{\rot}[1]{\textcolor{red}{#1}}
\newcommand{\orange}[1]{\textcolor{orange}{#1}}
\newcommand{\gelb}[1]{\textcolor{yellow}{#1}}
\newcommand{\gruen}[1]{\textcolor{green}{#1}}
\newcommand{\cyan}[1]{\textcolor{cyan}{#1}}
\newcommand{\braun}[1]{\textcolor{brown}{#1}}
\newcommand{\violett}[1]{\textcolor{violet}{#1}}
\newcommand{\schwarz}[1]{\textcolor{black}{#1}}
\newcommand*\centermathcell[1]{\omit\hfil$\displaystyle#1$\hfil\ignorespaces}
\NewDocumentCommand{\dx}{o}{\,\mathrm{d}\IfNoValueTF{#1}{x}{#1}}
\renewcommand{\labelenumi}{({\roman{enumi}})}
\renewcommand{\theenumi}{({\roman{enumi}})}
\renewcommand{\labelitemi}{\textendash}
\newcommand{\dotted}[1]{%
	\tikz[baseline=(todotted.base)]{
		\node[inner sep=0pt,outer sep=2pt] (todotted) {#1};
		\draw[dotted] (todotted.south west) -- (todotted.south east);
	}%
}%

\newcommand\Overline[2][1pt]{%
	\begin{tikzpicture}[baseline=(a.base)]
		\node[inner xsep=0pt,inner ysep=2pt] (a) {$#2$};
		\draw[line width= #1] (a.north west) -- (a.north east);
	\end{tikzpicture}
}

\NewDocumentCommand{\qrRand}{O{1.75cm} m O{} O{0pt}}{%
	\setlength\marginparsep{#4}%
	\marginnote{%
		\centering
		\href{#2}{%
			\IfSubStr{#2}{geogebra}{%
				\fancyqr[color=black,height=#1,
				image={\scalebox{1.5}{\vectorlogo[icon]{geogebra}}},
				image padding=0]{#2}%
			}{%
				\fancyqr[color=black,height=#1,image padding=0]{#2}%
			}%
		}%
		\ifstrempty{#3}{}{%
			\\[2pt]%
			\begin{tikzpicture}
				\node[draw,thick,rounded corners=10pt,
				minimum width=#1,minimum height=.75cm,
				inner sep=2pt,align=center]{%
					\adjustbox{max width=\dimexpr#1-4pt\relax,
						max height=\dimexpr.75cm-4pt\relax}{\texttt{#3}}%
				};
			\end{tikzpicture}%
		}%
	}%
}

\makeatletter
\newcommand{\discardpages}[1]{%
	\xdef\discard@pages{\zap@space#1 \@empty}%
	\AtBeginShipout{%
		\@expandtwoargs\in@{,\the\c@page,}{,\discard@pages,}%
		\ifin@\AtBeginShipoutDiscard\fi
	}%
}
\newcommand{\keeppages}[1]{%
	\xdef\keep@pages{\zap@space#1 \@empty}%
	\AtBeginShipout{%
		\@expandtwoargs\in@{,\the\c@page,}{,\keep@pages,}%
		\ifin@\else\AtBeginShipoutDiscard\fi
	}%
}
\def\ifGm@preamble#1{\@firstofone}
\appto\restoregeometry{%
	\ifLuaTeX
		\pagewidth=\paperwidth
		\pageheight=\paperheight
	\else
		\pdfpagewidth=\paperwidth
		\pdfpageheight=\paperheight
	\fi}
\apptocmd\newgeometry{%
	\ifLuaTeX
		\pagewidth=\paperwidth
		\pageheight=\paperheight
	\else
		\pdfpagewidth=\paperwidth
		\pdfpageheight=\paperheight
	\fi}{}{}
\makeatother

\newcolumntype{M}[1]{>{\centering\arraybackslash}m{#1}}

\makeatletter
\newcommand{\LastPageNumber}{%
	\zref@extractdefault{LastPage}{abspage}{1}%
}
\makeatother

\newcommand{\seitenzahl}{~%
	\ifnum\LastPageNumber=1\relax
	\else
	(\thepage/\LastPageNumber)%
	\fi
}
\newcommand{\hilfe}{\textit{Gar keine Idee? $\to$ Hilfekarte}}
\newcounter{tippcounter}
\newcounter{tippcounterdone}
\newcommand{\tipp}{\stepcounter{tippcounter}\textit{Hilfe? $\to$ Tipp\thetippcounter~unten}}
\newcommand{\tipptext}[1]{\vspace*{\fill}\hrule\vspace*{0.1cm}
	\begin{sideways}\begin{sideways}
			\begin{minipage}{\textwidth}\footnotesize\stepcounter{tippcounterdone}		
				\textit{Tipp\thetippcounterdone:}~#1
			\end{minipage}
	\end{sideways}\end{sideways}
}
\newboolean{loesung}
\setboolean{loesung}{false}

\newdimen\lsgHoehe
\newdimen\lsgBreite
\newdimen\lsgVorhoehe
\newbox\lsgVorbox

\makeatletter
\NewDocumentEnvironment{lsg}{O{0cm} O{0pt} g +b}{%
	\let\lsg\lsg@short
	\par
	\lsgVorhoehe=0pt
	\IfNoValueTF{#3}{%
		\ifx\height\fb@frh
			\endgroup\egroup
			\lsgVorhoehe=\dimexpr\ht\@tempboxa+\dp\@tempboxa+\height\relax
			\setbox\lsgVorbox=\box\@tempboxa
			\setbox\@tempboxa\vbox\bgroup
			\advance\hsize-\width\FrameRestore
			\textwidth\hsize
			\columnwidth\hsize
			\unvbox\lsgVorbox
			\begingroup\def\@currenvir{lsg}%
		\fi
		\lsgBreite=\dimexpr\linewidth-#1\relax
		\ifdim\pagegoal=\maxdimen
			\lsgHoehe=\dimexpr\textheight-#2-\lsgVorhoehe\relax
		\else
			\lsgHoehe=\dimexpr\pagegoal-\pagetotal-#2-\lsgVorhoehe\relax
		\fi
		\@tempdima=\dimexpr\textheight-#2-\lsgVorhoehe\relax
		\ifdim\lsgHoehe>\@tempdima\lsgHoehe=\@tempdima\fi
		\ifdim\lsgHoehe<0pt\lsgHoehe=0pt\fi
	}{%
		\lsgBreite=\dimexpr\linewidth-#1\relax
		\lsgHoehe=#3\relax
	}%
	\noindent
	\hspace*{#1}%
	\parbox[t][\lsgHoehe][t]{\lsgBreite}{%
		\ifthenelse{\boolean{loesung}}{%
			\settasks{label-format=\color{red},item-format=\color{red}}%
			\color{red}#4%
		}{%
			\strut
		}%
		\vfill
	}%
	\par
}{}
\let\lsg@environment\lsg
\NewDocumentCommand{\lsg@short}{m O{}}{%
	\ifthenelse{\boolean{loesung}}%
	{\textcolor{red}{#1}}%
	{#2}%
}
\let\lsg\lsg@short
\AddToHook{env/lsg/begin}{\let\lsg\lsg@environment}
\makeatother

\newcommand\colorxcancel[2][black]{\renewcommand\CancelColor{\color{#1}}\xcancel{#2}}
\newcommand{\janein}[2]{\begin{tikzpicture}[baseline]\bfseries
		\draw (0,0) rectangle node[pos=.5] {\scriptsize \ifthenelse{\boolean{loesung}}{\ifthenelse{\equal{#2}{W}}{\colorxcancel[red]{\textcolor{black}{W}}}{\textcolor{black}{W}}}{\textcolor{black}{W}}} (#1,#1);
		\draw (#1+0.1,0) rectangle node[pos=.5] {\scriptsize \ifthenelse{\boolean{loesung}}{\ifthenelse{\equal{#2}{F}}{\colorxcancel[red]{\textcolor{black}{F}}}{\textcolor{black}{F}}}{\textcolor{black}{F}}} (2*#1+0.1,#1);
		\end{tikzpicture}}

\newif\ifintask
\newcount\liniertcount
\newcount\liniertFrageZeilen
\newdimen\liniertRest
\newdimen\liniertH
\newdimen\liniertTaskRest

\ExplSyntaxOn
\cs_new_eq:NN \liniertTaskAlt:nnn \__tasks_task:nnn
\cs_set_protected:Npn \__tasks_task:nnn #1#2#3
{
    \global\liniertTaskRest=
        \ifdim\pagegoal=\maxdimen
            \textheight
        \else
            \dimexpr\pagegoal-\pagetotal\relax
        \fi
    \intasktrue
    \liniertTaskAlt:nnn{#1}{#2}{#3}
    \intaskfalse
}
\ExplSyntaxOff

\NewDocumentCommand{\liniert}{o +g}{%
    \begingroup
    \par
    \liniertFrageZeilen=\prevgraf
    \vspace{.9\baselineskip}%
    \noindent
    \IfNoValueTF{#1}{%
        \liniertRest=
        \ifintask
            \liniertTaskRest
        \else
            \ifdim\pagegoal=\maxdimen
                \textheight
            \else
                \dimexpr\pagegoal-\pagetotal\relax
            \fi
        \fi
        \ifintask
            \dimen0=\baselineskip
            \multiply\dimen0 by\liniertFrageZeilen
            \advance\liniertRest by-\dimen0
            \advance\liniertRest by-.9\baselineskip
        \fi
        \liniertcount=0
        \ifdim\liniertRest>\baselineskip
            \advance\liniertcount by1
            \advance\liniertRest by-\baselineskip
            \loop
                \ifdim\liniertRest>1.9\baselineskip
                    \advance\liniertcount by1
                    \advance\liniertRest by-1.9\baselineskip
            \repeat
        \fi
    }{%
        \liniertcount=#1\relax
    }%
    \ifnum\liniertcount<1
        \liniertcount=1
    \fi
    \liniertH=1.9\baselineskip
    \multiply\liniertH by\liniertcount
    \def\L{%
        \loop
            \ifnum\liniertcount>0
                \rule{\linewidth}{.6pt}%
                \par
                \vspace{.9\baselineskip}%
                \advance\liniertcount by-1
        \repeat
    }%
    \IfNoValueTF{#2}{%
        \L
    }{%
        \ifthenelse{\boolean{loesung}}{%
            \setlength{\fboxsep}{8pt}%
            \noindent
            \fbox{%
                \parbox[t]
                    [\dimexpr\liniertH-2\fboxsep-2\fboxrule\relax]
                    [t]
                    {\dimexpr\linewidth-2\fboxsep-2\fboxrule\relax}
                    {%
                        \settasks{
                            label-format=\color{red},
                            item-format=\color{red}
                        }%
                        \color{red}#2%
                    }%
            }%
            \par
        }{%
            \L
        }%
    }%
    \endgroup
}

\newlength{\laengeluecke}
\newlength{\lueckebreite}
\newcounter{lueckencounter}
\newcommand{\luecke}[1]{%
	\begingroup\stepcounter{lueckencounter}%
	\setbox0=\hbox{\strut #1}\setlength{\laengeluecke}{\wd0}%
	\setlength{\lueckebreite}{2\laengeluecke}\ifdim\lueckebreite<1cm\setlength{\lueckebreite}{1cm}\fi
	\def\LueckeInhalt{\dotted{\ifthenelse{\boolean{loesung}}{%
				\makebox[\lueckebreite][c]{\strut{\color{red}#1}}%
			}{%
				\TextField[bordercolor={},backgroundcolor={},color=black,align=1,width=\lueckebreite,name=Luecke\thelueckencounter,value={}]{}%
	}}}%
	\ifmmode\hbox{\LueckeInhalt}\else\LueckeInhalt\fi\endgroup
}

\newcommand{\squareX}[1]{%
	\makebox[0pt][l]{\raisebox{.05cm}{\ensuremath{%
				\ifthenelse{\boolean{loesung}\AND\equal{#1}{W}}{\color{red}\boldsymbol{\times}}{}%
	}}}%
	\scalebox{1.15}{\ensuremath{\square}}%
}

\def\vektor#1{\left(\vcenter{\halign{\hfil$##$\hfil\cr \vektorA#1,,}}\right)}
\def\vektorA#1,{\if,#1,\else #1\cr \expandafter \vektorA \fi}
\ExplSyntaxOn
\NewDocumentCommand\punkt{m}{
	\left(
	\clist_use:nn{\;#1\;}{\;\middle|\;}
	\right)}
\ExplSyntaxOff

\tikzset{schnipselcard/.style={draw=black!35,fill={rgb,255:red,255;green,249;blue,196},rounded corners=2mm,line width=.4pt,blur shadow,inner sep=4mm,outer sep=0pt}}
\ExplSyntaxOn
\seq_new:N\g__s_seq \box_new:N\l__s_i \box_new:N\l__s_r
\dim_new:N\l__s_w \dim_new:N\l__s_g \dim_new:N\l__s_gap
\NewDocumentCommand{\schnipselzeile}{+m}{\seq_gput_right:Nn\g__s_seq{#1}}
\cs_new_protected:Npn\__s_mk:n#1{\ifthenelse{\boolean{loesung}}{\hbox_set:Nn\l__s_i{\tikz[baseline]{\node[schnipselcard]{\begingroup\def\[{ \(\displaystyle }\def\]{ \)}\begin{varwidth}{.85\linewidth}\centering#1\end{varwidth}\endgroup};}}}{\pgfmathrandominteger{\A}{-20}{20}\hbox_set:Nn\l__s_i{\rotatebox[origin=c]{\A}{\tikz[baseline]{\node[schnipselcard]{\begingroup\def\[{ \(\displaystyle }\def\]{ \)}\begin{varwidth}{.85\linewidth}\centering#1\end{varwidth}\endgroup};}}}}}
\cs_new_protected:Npn\__s_flush:{\dim_compare:nNnF{\l__s_w}={0pt}{\par\noindent\hbox to\linewidth{\hss\box_use:N\l__s_r\hss}\par\vskip\l__s_gap\hbox_set:Nn\l__s_r{}\dim_zero:N\l__s_w}}
\NewDocumentEnvironment{schnipsel}{O{12345} O{1}}
{\seq_gclear:N\g__s_seq}
{\pgfmathsetseed{#1}\sys_gset_rand_seed:n{#1}\dim_set:Nn\l__s_gap{\fp_to_dim:n{4*(#2)}}\hbox_set:Nn\l__s_r{}\dim_zero:N\l__s_w
	\ifthenelse{\boolean{loesung}}
	{\seq_map_inline:Nn\g__s_seq{\__s_mk:n{##1}\par\noindent\hbox to\linewidth{\hss\box_use:N\l__s_i\hss}\par\vskip\l__s_gap}}
	{\seq_gshuffle:N\g__s_seq\seq_map_inline:Nn\g__s_seq{\__s_mk:n{##1}\dim_set:Nn\l_tmpa_dim{\box_wd:N\l__s_i}\dim_set:Nn\l__s_g{\dim_compare:nNnTF{\l__s_w}={0pt}{0pt}{\l__s_gap}}\dim_compare:nNnT{\l__s_w+\l__s_g+\l_tmpa_dim}>\linewidth{\__s_flush:\dim_zero:N\l__s_g}\hbox_set:Nn\l__s_r{\tex_unhbox:D\l__s_r\kern\l__s_g\vtop{\box_use:N\l__s_i}}\dim_set:Nn\l__s_w{\l__s_w+\l__s_g+\l_tmpa_dim}}\__s_flush:}}
\ExplSyntaxOff

\NewDocumentCommand{\hilfsmittelfrei}{O{1}}{%
	\ensuremath{\vcenter{\hbox{\hilfsmittelfreiTikz[#1]}}}%
}
\NewDocumentCommand{\hilfsmittelfreiTikz}{O{1}}{%
	\tikz[x=1ex,y=1ex,scale=#1,line cap=round,line join=round]{%
		\draw[gray,line width=0.08ex,rounded corners=0.25ex](0,0)rectangle(1.8,2.8);
		\draw[gray,line width=0.06ex,rounded corners=0.15ex](0.25,2.05)rectangle(1.55,2.55);
		\foreach\x in{0.45,0.90,1.35}{\foreach\y in{0.45,0.90,1.35}{\fill[gray](\x,\y)circle(0.10);}}
		\draw[gray,line width=0.18ex](-0.25,-0.25)--(2.05,3.05)(-0.25,3.05)--(2.05,-0.25);
	}%
}

\NewDocumentCommand{\differenzierung}{m}{%
	\tikz[baseline=-0.6ex]{
		\ifcase#1
		\or
		\draw (0,0) circle (0.8ex);
		\or
		\draw (0,0) circle (0.8ex);
		\fill (0,0) -- (180:0.8ex) arc (180:360:0.8ex) -- cycle;
		\or
		\fill (0,0) circle (0.8ex);
		\fi
	}%
}

\newcounter{besumme}\newcommand{\BE}[1]{\IfSubStr{#1}{BE:}{\StrBehind{#1}{BE:}[\x]\IfSubStr{\x}{,}{\StrBefore{\x}{,}[\x]}{}\addtocounter{besumme}{\numexpr\x\relax}}{}}\makeatletter\newcommand{\BESumme}{\@ifundefined{gespeicherteBESumme}{0}{\gespeicherteBESumme}}
\AtEndDocument{\immediate\write\@auxout{\string\gdef\string\gespeicherteBESumme{\arabic{besumme}}}}\makeatother
\let\AufgabeAlt\Aufgabe\let\endAufgabeAlt\endAufgabe\RenewDocumentEnvironment{Aufgabe}{O{}}{\BE{#1}\AufgabeAlt[#1]}{\endAufgabeAlt}

\Form
